\ficha{Crise, 1913}

\bloco{Fontes lidas}

Não há volume parlamentar para 1913. Li o relato da reunião da Associação
Comercial do Rio de Janeiro publicado no \textit{Correio Paulistano} de 10 de
agosto de 1913 e as duas versões do debate no Senado em dezembro: o resumo de
\textit{O Paiz} e a transcrição do \textit{Correio da Manhã}, ambos de 28 de
dezembro de 1913.

\bloco{Síntese}

Com as retiradas de ouro e a crise do Tesouro, o comércio do Rio discute em
agosto se pede ao governo uma emissão de papel-moeda. As posições antecipam as
de 1914. Em dezembro, no Senado, a criação da Caixa é rediscutida em
retrospecto, e os chefes políticos que a fizeram divergem sobre qual teria sido
a taxa certa em 1906.

\bloco{A reunião da Associação Comercial, agosto}

O ministro da Fazenda, Rivadavia Corrêa, não fala na reunião. É o barão de
Ibirocahy quem relata que o ministro lhe declarara que certos depósitos não
podiam ser utilizados, porque isso seria \tr[I:per090972_1913_17979:2]{faltar á
fé dos contractos} (\jor{cp19130810}{Correio Paulistano}{10 ago. 1913}{2}). Os
depósitos em causa são os de Londres, que Miguel Calmon diz estarem à
disposição do governo; não é certo que se trate do ouro da Caixa.

Carlos Peixoto \tr[I:per090972_1913_17979:2]{vota contra qualquer emissão} de
papel-moeda. Miguel Calmon o acompanha, por razão que envolve a Caixa:

\cit[I:per090972_1913_17979:2]{uma emissão de papel contraria a organização da
Caixa de Conversão}{\jor{cp19130810}{Correio Paulistano}{10 ago. 1913}{2}}

Augusto Ramos, autor de artigo sobre a crise no \textit{Jornal do Commercio},
defende a emissão sem abrir mão da estabilidade:

\cit[I:per090972_1913_17979:2]{não é contrario ao cambio estavel como está.
Continua sendo favoravel a emissão de papel}{\jor{cp19130810}{Correio
Paulistano}{10 ago. 1913}{2}}

\bloco{O retrospecto no Senado, dezembro}

No debate em que Pinheiro Machado reconta a criação da Caixa (ficha 3), três
senadores divergem sobre a taxa de 1906. Francisco Glycério diz que a oposição
ao projeto impediu que se adotasse \tr[I:per089842_1913_05446:3]{a taxa
verdadeira, que era a de 12}. Alcindo Guanabara apoia e insiste:

\cit[I:per178691_1913_10674:2]{Havia muito mais probabilidade de ser 12 do que
de 15}{\jor{opaiz19131228}{O Paiz}{28 dez. 1913}{2}}

Pinheiro Machado responde que as mesmas resistências impediram que se
mantivesse, sob Nilo Peçanha, \tr[I:per089842_1913_05446:3]{a taxa de 15, que
era a taxa verdadeira} (\jor{cm19131228}{Correio da Manhã}{28 dez. 1913}{3}).

\bloco{Achados}

\begin{enumerate}[leftmargin=*,nosep]
\item O desenho do debate de 1914 já está posto em agosto de 1913: emissão
contra a organização da Caixa, emissão com câmbio estável, recusa de qualquer
emissão.
\item A frase de Rivadavia sobre a fé dos contratos é relato de terceiro e se
refere a depósitos em Londres. Não serve, sem outra fonte, como posição do
ministro sobre o troco das notas da Caixa.
\item Em 1913 os autores políticos da Caixa ainda discutem se a taxa de 1906
deveria ter sido 12 ou 15. Alcindo mantém a posição de 1906.
\end{enumerate}

\bloco{Limites}

As duas versões do debate de dezembro diferem na forma: \textit{O Paiz} resume
em terceira pessoa, o \textit{Correio da Manhã} transcreve em primeira, e a
página deste está mutilada no fim. A reunião de agosto é conhecida por telegrama
do Rio publicado em São Paulo.
